\section{Related Work}
\label{sec:related-work}

\subsection{Task Completion and Program Understanding}

Many benchmarks evaluate model capabilities through the completion of programming tasks. HumanEval, MBPP, and APPS check whether generated code passes tests~\cite{chen2021codex,austin2021program,hendrycks2021apps}, while CodeXGLUE covers both understanding and generation~\cite{lu2021codexglue}. LiveCodeBench also includes execution and test-output prediction~\cite{jain2025livecodebench}, and BigCodeBench emphasizes the composition of library calls~\cite{zhuo2025bigcodebench}. At the repository level, SWE-bench evaluates issue resolution~\cite{jimenez2024swebench}, while Defects4J provides real Java defects and tests for repair studies~\cite{just2014defects4j}. \benchmark instead examines executable relationships between methods and the evidence returned with each reachability decision.

\subsection{Runtime Behavior and Dependencies}

Other benchmarks focus on specific properties of program behavior. CRUXEval asks models to predict inputs or outputs for short Python functions, while REval examines predictions about runtime behavior and intermediate states~\cite{gu2024cruxeval,chen2025reval}. Xie et al.'s CoRE evaluates static data and control dependencies and information flow, whereas Gao et al.'s CoRE tests implementation invariance and intermediate execution-state prediction~\cite{xie2025core,gao2026core}. SemBench includes function reachability through connectivity in a static call graph over real-world C programs~\cite{xu2026semb}. TraceEval asks models to recover runtime call structures from repository-derived programs, whose edges are verified through execution~\cite{li2026traceeval}. CABRA also controls task difficulty through generated call-graph transformations, but evaluates code-editing tasks with coding agents~\cite{balepur2026cabra}. Table~\ref{tab:benchmark-positioning} summarizes these task distinctions.

\begin{table}[t]
\centering
\caption{Task scope of benchmarks related to reachability and execution reasoning. Each row describes the question evaluated and the programs and oracles used.}
\label{tab:benchmark-positioning}
\scriptsize
\setlength{\tabcolsep}{3pt}
\begin{tabularx}{\columnwidth}{@{}p{0.23\columnwidth}p{0.37\columnwidth}X@{}}
\toprule
Benchmark & Question & Program and oracle \\
\midrule
SemBench~\cite{xu2026semb} & Is a function reachable in a static call graph? & Real-world C programs with static semantic labels. \\
CoRE (Xie et al.)~\cite{xie2025core} & Are statements related by data or control dependencies? & Human-verified static-analysis tasks. \\
TraceEval~\cite{li2026traceeval} & What is the program's runtime call structure? & Repository-derived programs with call edges verified by execution. \\
\benchmark & Does a target execute from a source under the values of a generated execution? & Parameter-controlled Java classes with execution-validated labels and graphs. Answers and traces are scored separately. \\
\bottomrule
\end{tabularx}
\end{table}


\benchmark asks whether a target method executes from a source under values supplied by one generated execution, rather than evaluating static reachability or recovering the complete runtime call structure. Its construction parameters control class size, guarded execution, and values passed between methods, which allows us to compare model behavior across structural difficulties. The Boolean answer and method trace are evaluated separately, so we can check whether a correct answer is supported by methods connected through executable calls. However, the generated classes cover a narrower range of Java behavior than repository-derived corpora.

\subsection{Reachability and Call-Graph Analysis}

Call-graph analysis studies the relationships used to determine possible calls~\cite{grove2001framework,tip2000scalable}. In our tasks, these relationships concern calls that can execute under a fixed invocation context. The graph representation therefore allows us to check the model's exploration, including whether it follows visible calls disabled by conditions or omits methods reachable from the source.
